\part{Глава №2. Основы архитектуры, операционных систем и инструментов программиста.}

\chapter{Тема №3. Основы архитектуры компьютера.}
\section{Лекция №4. Как устроен компьютер и где исполняется программа.}

\begingroup
\setlength{\emergencystretch}{2em}
\label{lecture04}
\subsection*{Цели и маршрут занятия (90 минут)}
В первых трёх лекциях мы изучали, как биты представляют данные и результаты
вычислений. Теперь соединим эти знания с устройством компьютера: кто выполняет
операции, где находятся команды и данные и как они перемещаются между
компонентами. Источник программы занятия --- раздел 2 ПУД курса.

После лекции студент сможет объяснить назначение CPU, RAM, накопителей и GPU;
различить адрес и содержимое памяти; проследить чтение элемента массива;
объяснить роль кэша, контроллера, DMA и прерывания; прочитать короткую
последовательность инструкций NASM. Учебная платформа для конкретных примеров
--- Linux x86-64 с моделью типов LP64. Условные адреса на рисунках не являются
адресами, которые следует подставлять в запускаемую программу.

\begin{center}
\small
\begin{tabular}{|p{0.73\linewidth}|r|}
\hline
Блок & Минуты \\
\hline
Общая картина: инструкции и данные & 0--3 \\
Компоненты компьютера & 3--15 \\
CPU и архитектурные модели & 15--23 \\
Шины и контроллеры & 23--33 \\
Иерархия памяти и кэш & 33--45 \\
Ввод-вывод, DMA, прерывания & 45--54 \\
Адреса, типы, указатели, массивы & 54--62 \\
Сквозной пример C и NASM & 62--76 \\
Стек, куча, время жизни & 76--82 \\
Система целиком & 82--87 \\
Итоги и самопроверка & 87--90 \\
\hline
\end{tabular}
\end{center}

\subsection{Общая картина: два потока}
Чтобы выполнить \texttt{a[1] += 5}, процессору нужны и \emph{инструкция},
описывающая действие, и \emph{данные}, над которыми выполняется действие.
Машинные инструкции тоже представлены байтами. Они не возникают внутри
CPU сами собой: подсистема выборки получает их из памяти.

\begin{quote}
Программа хранится на накопителе; при запуске ОС организует её размещение
в адресном пространстве процесса. Инструкции исполняет процессор, используя
регистры, кэши и оперативную память. Устройства выполняют ввод-вывод.
\end{quote}

Не следует смешивать \textbf{исходный текст} C или NASM, \textbf{исполняемый
файл} с машинными инструкциями и \textbf{процесс} --- выполняющийся экземпляр
программы со своим состоянием. Путь от текста к машинным инструкциям будет
подробно разобран на лекции №7, роль ОС при запуске --- на лекции №5.

\subsection{Компоненты компьютера}
\subsubsection*{CPU и RAM}
\textbf{CPU} (Central Processing Unit, центральный процессор) выполняет
машинные инструкции: арифметику, сравнения, переходы, обращения к памяти.
В одном процессоре может быть несколько \textbf{ядер}. Каждое ядро имеет
средства исполнения инструкций; часть ресурсов, например кэш последнего
уровня, может быть общей. Ядро, логический процессор и физический корпус
процессора --- разные понятия. SMT позволяет одному ядру поддерживать
несколько аппаратных потоков, но не удваивает все исполнительные ресурсы.

\textbf{RAM} (Random Access Memory) --- оперативная память. В обычном ПК это
DRAM: она хранит команды и данные, нужные работающим программам, и требует
питания и периодического обновления заряда. При выключении питания содержимое
обычной RAM теряется. Больший объём RAM позволяет держать больше рабочих
данных, но сам по себе не ускоряет каждое сложение в процессоре.

\textbf{Random access} означает возможность обращаться по адресу, а не
обязательный случайный порядок обращений. Время доступа к реальной DRAM
зависит от состояния контроллера, банков и характера запросов.

\subsubsection*{HDD: механика имеет значение}
\textbf{HDD} хранит данные на вращающихся магнитных пластинах. Для чтения
нужного участка головку требуется переместить к дорожке и дождаться, пока
нужный сектор окажется под ней. Поэтому мелкие случайные обращения дороги:
механическое ожидание измеряется миллисекундами. При последовательном чтении
большого участка эти затраты распределяются на много байтов.

Например, при 7200 оборотах в минуту один оборот занимает
$60/7200\approx0.00833$ с, то есть 8.33 мс. В упрощённой модели среднее
ожидание нужного положения равно половине оборота, около 4.17 мс.
Это \emph{только вращательная задержка}: поиск дорожки, передача и очередь
запросов добавляют время. Контроллер и кэш накопителя также влияют на результат.

\subsubsection*{SSD: флеш-память и контроллер}
\textbf{SSD} обычно использует NAND Flash. Механического перемещения нет,
поэтому случайный доступ значительно быстрее HDD. Однако SSD --- не RAM:
время доступа, интерфейс и правила записи иные. Данные программируются
страницами, а стираются более крупными блоками. Контроллер отображает
логические адреса блоков на физические области Flash, исправляет ошибки
и распределяет износ. Перезапись логического блока не обязательно означает
немедленную перезапись того же физического места.

Типичные характеристики накопителя: ёмкость, задержка, скорость
последовательной передачи и число операций ввода-вывода в секунду (IOPS).
Одно число «MB/s» не описывает все нагрузки. И HDD, и SSD сохраняют данные
без питания, но данные, ещё находящиеся в энергозависимых буферах, могут
требовать отдельного подтверждения сохранения.

\subsubsection*{GPU и видеопамять}
\textbf{GPU} (Graphics Processing Unit) оптимизирован для высокой суммарной
производительности при большом числе подходящих параллельных операций.
Например, над пикселями или элементами матрицы можно выполнять сходные
вычисления. CPU ориентирован в том числе на быстрый ход одного потока
со сложным управлением. Это различие приоритетов, а не запрет CPU на
параллелизм или GPU на условные переходы.

\textbf{Дискретный GPU} обычно имеет собственную видеопамять (VRAM).
Передача данных между RAM и VRAM и запуск работы имеют стоимость: маленькую
задачу может оказаться быстрее выполнить на CPU. \textbf{Встроенный GPU}
обычно использует общую системную RAM, хотя конкретная организация зависит
от платформы. Больший объём VRAM означает возможность хранить больше данных,
а не автоматически большую скорость вычисления.

Путь изображения в упрощённом виде: CPU и драйвер готовят команды и ресурсы;
GPU выполняет работу и формирует изображение в буфере; подсистема вывода
считывает его и передаёт сигнал дисплею. Дисплей не является памятью CPU,
а не вся работа GPU обязательно связана с показом изображения.

\subsubsection*{Плата, периферия, прошивка, питание}
Материнская плата обеспечивает соединения, разъёмы и питание компонентов.
\textbf{Контроллер} управляет обменом с памятью или устройством. В современных
ПК контроллер DRAM и часть линий PCIe часто находятся внутри CPU; чипсет
подключает дополнительные устройства. Конкретная топология различается.

Клавиатура, мышь, сетевой адаптер, накопитель и дисплей взаимодействуют
с системой через свои интерфейсы и контроллеры. \textbf{Драйвер} --- программный
компонент ОС, который знает, как работать с конкретным классом оборудования;
контроллер --- аппаратная часть.

Прошивка платформы, обычно реализующая интерфейсы \textbf{UEFI}, хранится
в энергонезависимой Flash, выполняет начальную инициализацию и помогает
передать управление загрузчику ОС. Исторический термин ROM означает память
только для чтения; современная память прошивки часто допускает обновление.
Блок питания обеспечивает нужные напряжения, охлаждение отводит тепло.
При перегреве процессор может снижать частоту: производительность зависит
и от условий работы, а не только от надписи на корпусе.

\subsection{Что находится внутри ядра CPU}
\begin{description}
\item[Регистры.] Небольшой набор непосредственно доступных инструкциям
хранилищ операндов, адресов и состояния. В x86-64 есть, например,
\texttt{RAX}, \texttt{RBX}, \texttt{RSP}; \texttt{EAX} обозначает младшие
32 бита \texttt{RAX}. Запись в \texttt{EAX} обнуляет старшие 32 бита
\texttt{RAX}. Регистры не являются обычными ячейками RAM с адресами C.
\item[Исполнительные устройства.] АЛУ выполняет целочисленные операции;
есть также устройства работы с плавающей точкой, векторными операциями,
адресами и памятью. Не каждая операция идёт через одно универсальное АЛУ.
\item[Управляющая логика.] Выбирает и декодирует инструкции, организует
их исполнение. Регистр \texttt{RIP} связан с адресом инструкции в
архитектурной модели, а \texttt{RFLAGS} содержит флаги состояния.
\item[Кэши.] Хранят копии инструкций и данных, сокращая обращения к RAM.
Разные уровни кэша имеют разный объём и стоимость доступа.
\end{description}

Учебный цикл: \textbf{выбрать инструкцию --- декодировать --- получить
операнды --- выполнить --- сохранить результат}. Переход меняет дальнейший
порядок исполнения. Это логическая модель: современный CPU использует
конвейер, предсказание переходов и перекрытие работы над несколькими
инструкциями. Одна инструкция не обязана занимать один такт.

При частоте 3 GHz такт длится примерно $1/(3\cdot10^9)$ с, или 0.333 нс.
Частота не равна числу выполненных инструкций в секунду. Важны число
инструкций задачи, степень параллелизма, пропуски кэша и ожидания.

\subsection{Общая и раздельная память команд и данных}
В \textbf{фон-неймановской модели} команды и данные хранятся в общей памяти
и используют общее адресное пространство. В \textbf{гарвардской модели}
память команд и данных, а также пути доступа к ним разделены. Один и тот же
числовой адрес в двух пространствах может обозначать разные места.
В строгой модели обычная команда чтения данных не читает память программ.
Реальные гарвардские процессоры могут предоставлять специальные инструкции
для доступа к коду как к данным.

\begin{center}
\small
\begin{tabular}{|p{0.28\linewidth}|p{0.58\linewidth}|}
\hline
Модель & Организация \\
\hline
Фон Нейман & Общая память команд и данных \\
Гарвард & Раздельные памяти и пути доступа \\
Модифицированная гарвардская & Разделение на части уровней при возможности общей памяти \\
\hline
\end{tabular}
\end{center}

Типичный современный x86-64 сочетает общее адресное пространство с
\emph{раздельными} L1-кэшами инструкций и данных. Поэтому классификация
зависит и от уровня, на котором рассматривается система.

\paragraph{Общая память не означает разрешение исполнять любой байт.}
В простой машине с общей памятью без защиты процессор может пытаться
интерпретировать байты доступной памяти как команды. В современной системе
права страниц ограничивают чтение, запись и исполнение. Например, бит
\texttt{NX/XD} запрещает выборку инструкций из страницы. При попытке
исполнения возникает исключение, а не нормальное выполнение этих байтов.
Стек и куча обычного процесса, как правило, неисполняемы. Точный набор
разрешений определяется архитектурой и ОС. Организация памяти и её защита
--- два разных вопроса.

\subsection{Шины, соединения и контроллеры}
\textbf{Шина} --- средство обмена между компонентами с согласованными
правилами передачи. Классическая учебная схема показывает CPU, RAM и
устройства на общей шине. В современном ПК несколько разных соединений,
в том числе выделенные каналы памяти и коммутируемые последовательные связи.

В логической модели обмена выделяют:
\begin{itemize}
\item \textbf{адрес}: к какому месту или ресурсу обращаемся;
\item \textbf{данные}: что передаём;
\item \textbf{управление}: чтение или запись, готовность, завершение, ошибки.
\end{itemize}
Это не обязательно три отдельных набора проводов: в пакетном соединении
адрес, данные и признаки операции могут передаваться по одним линиям.
Ширина адреса и ширина одновременно передаваемых данных тоже не обязаны
совпадать. Из «64-битный CPU» нельзя вывести, что у каждого его интерфейса
64 физических линии.

\paragraph{Чтение и запись.}
При чтении инициатор формирует запрос с адресом; адресат возвращает данные
и статус. При записи инициатор передаёт адрес и данные, а протокол задаёт,
как сообщить о принятии операции. CPU общается с DRAM через контроллер
памяти: тот переводит запросы в команды конкретных каналов и банков.
Наличие запроса ещё не означает немедленного получения результата.

\paragraph{Современная схема соединений.}
\begin{center}
\small
\begin{tabular}{|p{0.3\linewidth}|p{0.56\linewidth}|}
\hline
Соединение & Типичное назначение \\
\hline
Каналы DDR-памяти & Контроллер памяти CPU и модули DRAM \\
PCIe & GPU, NVMe SSD, сетевые адаптеры; линии и коммутаторы \\
SATA & Подключение HDD и SATA SSD через контроллер \\
USB & Периферия через хост-контроллер и, возможно, концентраторы \\
\hline
\end{tabular}
\end{center}
\textbf{M.2} --- спецификация форм-фактора и разъёма, а не синоним NVMe.
\textbf{NVMe} --- протокол работы с энергонезависимыми накопителями,
обычно через PCIe в локальном ПК. M.2-накопитель может использовать SATA
или PCIe/NVMe; поддержка зависит от устройства и слота. \textbf{PCIe}
не является только «шиной видеокарты»: его используют разные устройства.

\paragraph{Задержка и пропускная способность.}
Задержка --- время ожидания отдельного результата. Пропускная способность
--- объём передачи за единицу времени при заданной нагрузке.
Для грубой модели без перекрытия работы:
\[
T\approx L+\frac{V}{B},
\]
где $L$ --- начальная задержка, $V$ --- объём, $B$ --- скорость передачи.
Например, $L=100$ мкс, $V=1$ MB, $B=1$ GB/s дают около 1.1 мс.
Используем десятичные MB и GB; служебные передачи и очередь не учитываем.
Для маленького запроса доминирует задержка, для большого --- время передачи.
Много параллельных запросов могут повысить суммарную скорость, не уменьшая
задержку каждого. Это полезно помнить при сравнении CPU, GPU и накопителей.

\subsection{Иерархия памяти: скорость, объём и доступ}
\label{lecture04-memory}
Быструю память дорого делать большой. Поэтому компьютер сочетает несколько
уровней: небольшие регистры и кэши, большую RAM и ещё более вместительные
накопители. Эта иерархия --- компромисс, а не закон, что любая операция
на верхнем уровне всегда быстрее любой операции на нижнем.

\begin{center}
\small
\begin{tabular}{|p{0.19\linewidth}|p{0.27\linewidth}|p{0.4\linewidth}|}
\hline
Уровень & Типичная технология & Доступ и назначение \\
\hline
Регистры & Схемы внутри CPU & Указаны в инструкциях; текущие операнды \\
L1/L2/L3 & SRAM & Аппаратно управляемые копии линий памяти \\
RAM & DRAM & Адресуемая рабочая память программ \\
VRAM & Обычно GDDR/HBM & Рабочая память GPU; у встроенного GPU часто общая RAM \\
SSD/HDD & NAND Flash / магнитная запись & Блочный ввод-вывод, долговременное хранение \\
Прошивка & Обычно Flash & Начальная настройка и запуск системы \\
\hline
\end{tabular}
\end{center}

SRAM сохраняет состояние при наличии питания без периодического обновления
заряда, необходимого DRAM; обе технологии энергозависимы. «Статическая»
в названии SRAM не означает сохранение после выключения. Flash энергонезависима.
Регистры не следует механически называть ещё одним уровнем SRAM-кэша.

\paragraph{Ориентиры скорости, а не паспортные константы.}
Для современного настольного компьютера полезны следующие порядки величин:
типичные простые зависимые операции с регистрами и попадания в L1 ---
единицы тактов; L2 --- часто десятки тактов, L3 --- десятки и более;
обращение к DRAM --- десятки--сотни наносекунд; случайный доступ к SSD ---
десятки--сотни микросекунд, к HDD --- миллисекунды. Результат зависит от
процессора, устройства, нагрузки, очередей и попаданий в буферы.
Деление в регистрах, например, не обязано занимать один такт.

При 3 GHz ожидание 100 нс соответствует примерно 300 тактам.
CPU может перекрывать ожидание другой работой, но цепочка зависимых загрузок
часто не позволяет скрыть всю задержку. Доступность также означает
\emph{способ обращения}: инструкция загрузки из обычной памяти не заменяет
команду чтения сектора SSD. Отображённый в память файл может выглядеть для
программы как адресуемые байты, но его подгрузкой управляет ОС.

\subsection{Кэш процессора}
\textbf{Кэш} хранит копии блоков памяти, которые могут понадобиться скоро.
Попадание (cache hit) означает, что нужная линия найдена на проверяемом
уровне; промах (cache miss) требует дальнейшего поиска. Промах L1 не означает
автоматически обращение к RAM: данные могут находиться в L2, L3 или кэше
другого ядра в согласованной системе.

Типичная учебная схема: L1I для инструкций, L1D для данных, собственный L2
ядра и общий для нескольких ядер L3. Эта схема не универсальна. Она
объясняет, почему получение байтов команды и чтение её операнда --- разные
обращения. Программа обычно не распределяет свои переменные по L1/L2/L3
вручную: размещением управляет аппаратура.

\paragraph{Кэш-линия и локальность.}
Для примеров принимаем линию 64 байта --- типичный размер у современных
x86-64. Загрузка одного \texttt{int} может потребовать получения целой линии.
Если \texttt{sizeof(int)==4} и начало группы выровнено по границе линии,
в ней помещаются 16 элементов. Запрос 17 последовательных элементов
затрагивает уже две линии. Если начало смещено, даже 16 элементов могут
занять две линии.

\textbf{Пространственная локальность}: рядом с нужным адресом часто лежат
следующие нужные данные. \textbf{Временная локальность}: недавно использованное
значение может понадобиться снова. При последовательном обходе массива
полезно много байтов каждой линии. При шаге в 16 четырёхбайтовых элементов
каждое чтение попадает в новую линию, и из 64 байтов непосредственно
используются лишь 4. Это иллюстрация использования линий, а не обещание
ускорения ровно в 16 раз: влияют предвыборка, параллельные запросы и объём кэша.

\paragraph{Запись и dirty-линии.}
Для обычной кэшируемой памяти x86-64 распространена политика
\textbf{write-back}. Запись проходит через буфер записей; кэш-линия получает
новое значение и может стать \emph{грязной} (dirty), то есть отличаться от
копии на следующем уровне. Она будет передана дальше по правилам кэша,
например при вытеснении. Следующая загрузка на том же ядре может получить
собственную ещё ожидающую запись из буфера. Это не означает, что программа
произвольно получает старое значение после своего присваивания.

Правила видимости между ядрами сложнее и будут важны при многопоточности.
Главное сейчас: \textbf{запись в переменную, запись в DRAM и сохранение файла
на SSD --- не одно событие}. В нашем примере рассматривается одно ядро,
обычная память и отсутствие конкурирующих изменений.

\subsection{Ввод-вывод, DMA и прерывания}
\label{lecture04-interrupts}
Программа запрашивает действие у ОС. Драйвер настраивает контроллер,
передавая команду, параметры и сведения о буферах. Регистры устройства
могут быть отображены в адресное пространство (MMIO); это \emph{не}
обычные ячейки RAM и не регистры общего назначения CPU. На x86 есть
также отдельные инструкции портового ввода-вывода. Доступ к оборудованию
обычно выполняется привилегированным кодом ОС.

\paragraph{Опрос и уведомление.}
При \textbf{опросе} (polling) CPU периодически проверяет состояние устройства.
Это просто, но частые проверки расходуют процессорное время. При
\textbf{прерывании} (interrupt) система получает уведомление о событии.
Прерывания имеют накладные расходы, поэтому высокопроизводительный
ввод-вывод может сочетать их с опросом и обработкой пакетов событий.

\paragraph{DMA: передача без копирования каждого байта CPU.}
При прямом доступе к памяти (Direct Memory Access) устройство или
контроллер переносит блок в RAM либо из RAM после настройки со стороны
ОС. CPU не обязан выполнять отдельные \texttt{mov} для каждого байта.
Однако он участвует в подготовке операции и обработке результата.
Адрес для устройства не следует отождествлять с пользовательским указателем:
ОС организует допустимый буфер и необходимые преобразования.

Пример чтения, когда нужных данных ещё нет в файловом кэше ОС:
\begin{enumerate}
\item Программа запрашивает чтение; ОС и драйвер готовят запрос и буфер.
\item Контроллер накопителя читает данные и передаёт их в память по DMA.
\item Устройство сообщает о завершении, например прерыванием.
\item ОС отмечает завершение и делает результат доступным программе;
      в обычном буферизованном чтении возможна дополнительная копия.
\end{enumerate}
Если данные уже есть в файловом кэше RAM, физического чтения накопителя
может не потребоваться. Файловый кэш ОС и кэш CPU --- разные уровни.

\paragraph{Как обрабатывается аппаратное прерывание.}
\begin{enumerate}
\item Устройство, контроллер или таймер формирует запрос.
\item Механизм доставки направляет его нужному CPU; маскируемый запрос
      может ожидать, пока процессор сможет его принять.
\item CPU сохраняет предусмотренную архитектурой часть состояния и
      передаёт управление обработчику. Программный код сохраняет
      дополнительные регистры, которые будет использовать.
\item Обработчик выясняет причину, подтверждает или снимает запрос,
      фиксирует результат. Длительную работу можно отложить.
\item При возврате состояние восстанавливается; возможно продолжение
      прерванного кода. ОС может также решить запланировать другой поток.
\end{enumerate}
\textbf{Не каждое прерывание вызывает переключение процесса.} Прерывание
таймера даёт ОС возможность получить управление, учитывать время и принимать
решения планировщика. Прерывание завершения ввода-вывода сообщает о событии,
а не обязательно переносит сами данные.

Для USB-клавиатуры хост-контроллер организует обмен по USB и может уведомить
CPU о результате. Нажатие клавиши не вызывает напрямую функцию
пользовательской программы. Термин USB interrupt transfer тоже не следует
понимать как отдельный провод от клавиши к ядру CPU.

\paragraph{Прерывание, исключение, системный вызов.}
Аппаратное внешнее прерывание асинхронно относительно текущей инструкции.
\textbf{Исключение} связано с её выполнением: например, целочисленное деление
на ноль или страничное исключение. Последнее может быть штатным поводом
подгрузить страницу, а не ошибкой программы. \textbf{Системный вызов} ---
намеренное обращение программы к ОС; на Linux x86-64 обычно используется
инструкция \texttt{syscall}. Это не уведомление от периферии. У разных
архитектур названия и механизмы различаются; здесь важно различать причины.

\subsection{Байты, адреса и объекты C}
В примерах байт содержит 8 бит. \textbf{Адрес} указывает место, а
\textbf{содержимое} --- хранимые там биты. \textbf{Адресное пространство}
--- множество адресов, в котором выполняются обращения. Пользовательская
программа обычно работает с виртуальными адресами, переводимыми аппаратурой
при участии ОС в физические. Наличие 64-битного указателя не означает ни
наличие $2^{64}$ байтов RAM, ни допустимость всех его битовых комбинаций.

Машинное слово --- архитектурно обусловленная единица обработки; термин
нужно употреблять с указанием контекста. В NASM \texttt{word} всегда
16 бит, \texttt{dword} --- 32, \texttt{qword} --- 64. У x86-64 существуют
операции с разными размерами, а не только с 64 битами.

Для обычной Linux x86-64 LP64: \texttt{char} --- 1 байт, \texttt{int} --- 4,
\texttt{long} и указатель --- 8. Это свойства платформы. Например,
в распространённой 64-битной Windows используется LLP64 и \texttt{long}
обычно имеет размер 4 байта. Язык C гарантирует \texttt{sizeof(char)==1},
но число бит в байте определяется \texttt{CHAR\_BIT}.

\begin{lstlisting}[language=C,basicstyle=\ttfamily\small]
int x = 20;
int *p = &x;
*p = 25;
\end{lstlisting}
\texttt{x} --- объект типа \texttt{int}, \texttt{\&x} --- его адрес,
\texttt{p} --- переменная, хранящая указатель, \texttt{*p} --- обращение
к объекту по этому указателю. Присваивание \texttt{*p=25} меняет значение
\texttt{x}, но не адрес \texttt{x}. У самой переменной \texttt{p} также
есть собственное место хранения, если оно требуется реализации.

\paragraph{Массив.}
В \texttt{int a[3]=\{10,20,30\}} элементы идут подряд. Для элемента
с допустимым индексом $i$ байтовый адрес равен
\[
\operatorname{base}+i\cdot\operatorname{sizeof}(\text{элемента}).
\]
При условной базе \texttt{0x1000} и размере 4 байта адреса начала элементов
--- \texttt{0x1000}, \texttt{0x1004}, \texttt{0x1008}; размер массива ---
12 байтов. В выражении C \texttt{a+1} шаг уже масштабируется типом:
не нужно вручную умножать его ещё раз. Массив и указатель --- не один тип:
\texttt{sizeof(a)} здесь равно 12, а размер указателя на нашей платформе --- 8.

\paragraph{Порядок байт.}
Значение \texttt{0x12345678} типа \texttt{uint32\_t} на little-endian x86-64
занимает байты \texttt{78 56 34 12} по возрастающим адресам.
В big-endian было бы \texttt{12 34 56 78}. Меняется порядок байтов,
а не направление записи битов внутри каждого байта.
Для просмотра представления объекта в C можно использовать указатель
на \texttt{unsigned char}. Размеры печатают через \texttt{\%zu}, адреса ---
через \texttt{\%p} с аргументом, приведённым к \texttt{void *}.

\subsection{Сквозной пример: путь инструкций и данных}
\label{lecture04-trace}
Рассмотрим одно изменение массива:
\begin{lstlisting}[language=C,basicstyle=\ttfamily\small]
int a[3] = {10, 20, 30};
a[1] += 5;
\end{lstlisting}
Учебная реализация NASM, если \texttt{RBX} уже содержит адрес массива:
\begin{lstlisting}[basicstyle=\ttfamily\small]
mov eax, [rbx + 4]
add eax, 5
mov [rbx + 4], eax
\end{lstlisting}
Квадратные скобки обозначают обращение к памяти; без них \texttt{RBX}
обозначает регистр. Размер операнда памяти здесь выводится из \texttt{EAX}:
32 бита, то есть 4 байта. Это одна возможная реализация, а не обещание
буквального вывода компилятора. Оптимизатор может выбрать инструкцию
с операндом в памяти или вовсе вычислить известный результат заранее.

\subsubsection*{Шаг 0. От файла к доступным страницам}
Исполняемый файл находится на SSD/HDD. При запуске ОС создаёт отображения
кода и данных в виртуальном адресном пространстве. Весь файл не обязательно
сразу копируется в RAM: страницы могут подгружаться по требованию, а часть
данных создаётся при выполнении. В частности, локальный массив может быть
инициализирован инструкциями после входа в функцию.

Для дальнейших кадров предполагаем: страницы кода и данных уже доступны,
права позволяют исполнение кода и чтение/запись массива, исключений нет.
Условный адрес массива --- \texttt{0x1000}, условный адрес начала фрагмента
кода --- \texttt{0x400000}. Они виртуальные и выбраны только для рисунка.

\subsubsection*{Шаг 1. Процессор получает саму инструкцию}
В архитектурном разборе \texttt{RIP} задаёт адрес очередной инструкции.
Подсистема выборки получает байты кода через L1I. Если линия там отсутствует,
запрос обслуживается нижними уровнями и, при необходимости, RAM.
При холодном обращении удобно рисовать путь
\begin{center}
RAM $\longrightarrow$ нижние уровни кэша $\longrightarrow$ L1I
$\longrightarrow$ декодирование.
\end{center}
Не каждый реальный процессор копирует линию во все уровни именно в таком
порядке. Кэш другого ядра тоже может быть источником. На схеме показываем
смысл иерархии, а не универсальный протокол заполнения.

Для выбранного NASM-кодирования три инструкции имеют байты:
\begin{center}
\small
\begin{tabular}{|l|l|l|}
\hline
Условный адрес & Байты & Инструкция \\
\hline
\texttt{0x400000} & \texttt{8B 43 04} & \texttt{mov eax,[rbx+4]} \\
\texttt{0x400003} & \texttt{83 C0 05} & \texttt{add eax,5} \\
\texttt{0x400006} & \texttt{89 43 04} & \texttt{mov [rbx+4],eax} \\
\hline
\end{tabular}
\end{center}
Все три команды помещаются в одну 64-байтовую линию при данном выравнивании.
В этом примере длина каждой равна 3 байтам, но длина инструкций x86 переменная.
Аппаратура может заранее выбрать несколько команд; это не три обязательных
похода в RAM. Значение \texttt{05} уже содержится в коде инструкции сложения:
отдельная загрузка «переменной 5» из массива не нужна.

\subsubsection*{Шаг 2. Загрузка элемента}
\texttt{mov eax,[rbx+4]} вычисляет адрес \texttt{0x1004}. Перевод адреса
и проверка прав обеспечивают доступ к нужной памяти. Подсистема загрузок
ищет данные через L1D; при промахе --- дальше по иерархии.
В нашей учебной схеме при полном промахе линия данных поступает из RAM
через нижние уровни в L1D, а байты \texttt{14 00 00 00} дают число 20
в \texttt{EAX}. Остальные элементы массива не изменились.

\subsubsection*{Шаг 3. Сложение}
CPU декодирует \texttt{add eax,5}, получает значение регистра и непосредственный
операнд из инструкции. АЛУ вычисляет $20+5=25$. Результат попадает в
\texttt{EAX}; инструкция также обновляет арифметические флаги.
Здесь нет обращения к массиву за операндами сложения: 20 уже в регистре.
До следующей инструкции логическое значение \texttt{a[1]} ещё равно 20.

\subsubsection*{Шаг 4. Запись}
\texttt{mov [rbx+4],eax} направляет 4 байта результата в подсистему памяти.
После разрешённой записи программа должна наблюдать \texttt{a[1]==25}.
В обычной реализации запись проходит через буфер; линия L1D получает
байты \texttt{19 00 00 00}. При write-back копия в DRAM может некоторое
время содержать прежние байты: это не отдельное значение переменной,
которое обычная программа выбирает читать вместо согласованной подсистемы.
Файл на SSD от изменения массива не меняется.

\begin{center}
\small
\begin{tabular}{|l|l|l|}
\hline
Момент & \texttt{EAX} & Логическое содержимое массива \\
\hline
До фрагмента & Не используется & \texttt{10, 20, 30} \\
После загрузки & 20 & \texttt{10, 20, 30} \\
После сложения & 25 & \texttt{10, 20, 30} \\
После записи & 25 & \texttt{10, 25, 30} \\
\hline
\end{tabular}
\end{center}

\paragraph{Проверяемая демонстрация.}
Каталог \texttt{demos/lecture04/} содержит программу C, NASM-функцию
и Docker-окружение. Для проверки выполните из корня репозитория:
\begin{verbatim}
make lecture4-check
\end{verbatim}
Проверяются размеры типов, байты числа, размещение массива и результат
вызова NASM.
В вызываемой функции адрес поступает в \texttt{RDI} по System V AMD64 ABI;
обёртка сохраняет \texttt{RBX}, переносит туда адрес и восстанавливает регистр
перед возвратом. Эти дополнительные команды обеспечивают корректный вызов
из C; три центральные инструкции остаются ровно теми, что разобраны выше.
В демонстрации также проверяется их машинное кодирование.

\subsection{Стек и куча: где живут объекты}
\textbf{Стек} --- область, используемая реализацией для организации вызовов:
адресов возврата, сохранённых регистров, части локальных объектов.
В x86-64 \texttt{RSP} --- указатель стека. \texttt{call} сохраняет адрес
возврата в стеке, \texttt{ret} получает его оттуда. В System V AMD64 многие
аргументы передаются через регистры, поэтому «все аргументы на стеке» неверно.

Локальный объект с автоматическим временем хранения обычно существует
до выхода из блока. Реализация может держать его значение в регистре
или устранить объект: язык C не требует помещать каждую локальную переменную
на стек. Возвращать указатель на уже закончивший жизнь локальный объект
нельзя использовать как способ сохранить его данные.

\textbf{Кучей} обычно называют область и механизмы динамического выделения.
\texttt{malloc} предоставляет блок до явного \texttt{free}; объект не исчезает
только из-за выхода из функции, которая его выделила. Размещение и получение
памяти от ОС зависят от реализации аллокатора.
\begin{lstlisting}[language=C,basicstyle=\ttfamily\small]
int *p = malloc(3 * sizeof *p);
if (p != NULL) {
    p[0] = 10;
    p[1] = 20;
    p[2] = 30;
    p[1] += 5;
    free(p);
    p = NULL;
}
\end{lstlisting}
Нужен \texttt{<stdlib.h>}. Локальная переменная \texttt{p} и выделенный
блок --- разные объекты с разным временем жизни. После \texttt{free}
разыменование старого указателя недопустимо; присваивание \texttt{NULL}
этой переменной не исправляет другие копии указателя.
Глобальные объекты и статические локальные объекты имеют статическое время
хранения: деление «всё либо стек, либо куча» неполно.

\subsection{Компьютер целиком: чтение, вычисление, вывод}
Представим программу, которая читает массив чисел, изменяет их и рисует
диаграмму. Ввод-вывод загружает данные с накопителя в память; CPU получает
инструкции и данные через кэши. При подходящей большой задаче часть работы
можно передать GPU. Результат для отображения попадает в графическую
подсистему; для сохранения файла программа отдельно запрашивает запись у ОС.

Узкое место зависит от этапа: случайные чтения HDD ограничены механикой,
большая передача --- пропускной способностью, зависимые чтения RAM ---
задержкой, вычислительный цикл --- исполнительными ресурсами. Переход на SSD
не ускорит непосредственно сложение уже находящихся в регистрах операндов.
Больше RAM помогает при нехватке рабочей памяти; больше VRAM не заменяет
вычислительные блоки GPU. Оценивать нужно путь конкретной задачи.

\subsection*{Итоги и самопроверка}
\begin{enumerate}
\item Откуда берутся байты \texttt{mov eax,[rbx+4]}, а откуда загружаемое число?
\item Почему промах L1 не обязательно приводит к чтению DRAM?
\item Чем отличаются адрес \texttt{0x1004} и число 20?
\item Почему изменённый массив не сохраняется автоматически в файл?
\item Обязательно ли CPU копирует каждый байт при DMA?
\item Обязательно ли прерывание переключает процесс?
\item Почему общая память команд и данных не означает исполняемость стека?
\item Где находится объект, на который указывает локальная переменная \texttt{p}?
\end{enumerate}
\textbf{Краткие ответы:} код выбирается через подсистему инструкций, число ---
через подсистему данных; существуют нижние уровни кэша; адрес указывает место,
число является значением; нужна отдельная операция ввода-вывода; нет,
передачу выполняет устройство/контроллер; нет, это решение ОС; действуют
права страниц; зависит от того, на что указывает \texttt{p}, а не от
места хранения самого указателя.

\newpage
\subsection{Дополнительно: детали за пределами 90 минут}
\paragraph{Линия, индекс и тег.}
Для 64-байтовой линии младшие 6 бит байтового адреса задают смещение внутри
линии. В учебном кэше с прямым отображением и 64 наборами следующие 6 бит
выбирают набор, остальные образуют тег. Тег позволяет отличить разные
блоки, претендующие на одно место. Например, для адреса \texttt{0x1234}
смещение равно \texttt{0x34}=52, номер набора --- 8, тег --- 1.
Это модель кэша, а не описание конкретного L1 современного x86.

\paragraph{Ассоциативность и вытеснение.}
В наборно-ассоциативном кэше один набор может хранить несколько
линий с одним индексом. Английский термин --- set-associative cache.
Такая организация уменьшает некоторые конфликтные промахи.
При нехватке места выбирается линия для вытеснения. Политика не обязана
быть точным LRU. Три причины промаха: первое обращение к линии,
нехватка общего объёма и конкуренция за набор.

\paragraph{Политики записи.}
Write-through передаёт запись также на следующий уровень, write-back
позволяет отложить её до записи изменённой линии. Write-allocate описывает
получение линии при промахе записи; это другой вопрос. Обычная память,
MMIO и специальные инструкции могут иметь разные правила кэширования.

\paragraph{Согласованность кэшей.}
Несколько ядер могут иметь копии одной линии. Протокол согласованности
помогает организовать владение и распространение изменений. Даже изменения
разных переменных в одной линии могут вызывать лишние передачи
(false sharing). Согласованность кэшей не заменяет атомарные операции,
синхронизацию и модель памяти языка C.

\paragraph{Страницы и TLB.}
Виртуальное адресное пространство отображается на физическую память
страницами. В примере со страницей 4 KiB младшие 12 бит задают смещение,
старшие --- номер виртуальной страницы. Для \texttt{0x1234} это страница 1
и смещение \texttt{0x234}=564. TLB кэширует переводы адресов и сведения о
доступе, а не данные массива. Промах TLB может потребовать обхода таблиц;
страничное исключение возникает по иной причине, например при отсутствии
допустимого отображения. Размер страницы зависит от режима системы.

\paragraph{Выравнивание и структуры.}
У структуры \texttt{struct S \{ char c; int x; \};} при обычных правилах
выравнивания Linux x86-64 поле \texttt{x} имеет смещение 4, размер структуры
--- 8 байтов. Между полями есть padding. Для конкретной реализации это
проверяют через \texttt{offsetof} и \texttt{sizeof}; это не универсальный
результат для всех компиляторов и настроек упаковки.

\paragraph{Оптимизация и наблюдаемое поведение.}
Компилятор обязан сохранять требуемое языком наблюдаемое поведение, а не
по одной инструкции на каждую строку. Если результат массива нигде не
используется, его и вычисление могут исчезнуть. Код демонстрации вызывает
внешнюю NASM-функцию и проверяет результат; это делает разбор воспроизводимым.
\texttt{volatile} не следует использовать как универсальное средство
управления кэшем или синхронизации потоков.

\paragraph{Доставка событий.}
На современных x86 запросы могут доставляться через APIC, в том числе
сообщениями MSI/MSI-X от PCIe-устройств. ОС настраивает маршрутизацию и
обработчики. Эти детали не меняют базового различия между передачей данных
и уведомлением о готовности. Реальные системы используют очереди, пакетную
обработку и отложенную работу, чтобы уменьшать накладные расходы.

\paragraph{Источники.}
Общая модель компьютера и иерархия памяти: CS:APP~\cite{bryant-ohallaron-csapp}
и Patterson--Hennessy~\cite{patterson-hennessy}. Регистры, инструкции,
страницы и прерывания: Intel SDM~\cite{intel-manual}. Синтаксис и размеры
операндов: документация NASM~\cite{nasm-doc}. Объекты C, массивы и время
хранения: N1570~\cite{c11-draft}; соглашение вызова демонстрации:
System V AMD64 psABI~\cite{amd64-psabi}.

\endgroup
\begingroup
\setlength{\emergencystretch}{2em}
\subsection*{Архитектура и организация компьютера}
\label{lecture04-journal}
\textit{Сквозной пример: разработка журнала оценок группы.}
Напишем программу без гостевой ОС. Исполняемый код --- 16-битный NASM
в real mode на эмулированном QEMU PC (i440FX); C встречается лишь как
\emph{пояснение алгоритма}. Сохраняем действующий текст основной лекции №4;
здесь те же архитектурные понятия открываются при разработке шести
последовательных версий журнала. Студенты обозначены \texttt{S01}--\texttt{S16},
у каждого четыре работы и оценка 0--10 либо «не выставлена».

\begin{center}\small
\begin{tabular}{|p{0.23\linewidth}|p{0.59\linewidth}|}\hline
Минуты & Потребность и теория \\ \hline
0--7 & Задача и простая аппаратная модель \\
7--20 & Запуск: BIOS, сектор, CS:IP и первый символ \\
20--35 & Таблица: байты, адреса и текстовая видеопамять \\
35--48 & Ввод через BIOS и разбор скан-кода \\
48--60 & Собственный IRQ1 и очередь событий \\
60--72 & Средние, регистры, стек, код и кэш \\
72--86 & Долговременный сектор, BIOS-диск и проверка данных \\
86--90 & Итоги и связь с операционной системой \\ \hline
\end{tabular}
\end{center}

\paragraph{Контракт учебной машины.}
У QEMU PC 16 MiB RAM, BIOS, VGA-текстовый экран, PS/2-клавиатура,
PIC, COM1 для диагностики, загрузочная дискета 1.44 MiB и отдельный
IDE-совместимый диск 1 MiB для оценок. Настоящее устройство компьютера
сложнее, но выбранный набор позволяет увидеть, кто исполняет инструкции,
где лежат данные и как взаимодействуют компоненты. BIOS --- прошивка,
не гостевая ОС. Для тестов запускаем \texttt{qemu-system-i386}; название
команды не исключает исполнения 16-битных инструкций в real mode.
Ни одну характеристику производительности реального SSD или кэша не
измеряем по скорости QEMU.

\paragraph{Версия 1: начать исполнение и вывести одну оценку.}
После сброса CPU выбирает инструкции BIOS. BIOS выполняет базовую
инициализацию, выбирает дискету согласно порядку загрузки, размещает
её первый сектор по адресу \texttt{0000:7C00} и передаёт туда управление.
Его последние два байта имеют сигнатуру \texttt{55 AA}; сама программа
слишком велика для сектора 512 байтов. Поэтому наш загрузочный сектор
\texttt{boot16.asm} читает сектора 2--18 той же дорожки в RAM через
BIOS \texttt{INT 13h}, проверяет результат и делает дальний переход
к \texttt{1000:0000}:
\begin{lstlisting}[basicstyle=\ttfamily\small]
mov ax, 0x1000
mov es, ax
xor bx, bx
mov ah, 0x02
mov al, 17
mov cl, 2
int 0x13
jc boot_error
jmp 0x1000:0x0000
\end{lstlisting}
В полном листинге дополнительно задаются цилиндр, головка и номер
загрузочной дискеты \texttt{DL}, а также проверяется число прочитанных
секторов. За пределы первой дорожки не выходим; сборка отвергает программу
длиннее $17\cdot512=8704$ байтов. \textbf{Наш код запускается без GRUB},
хотя BIOS сам по себе остаётся частью учебной машины.

Адрес в real mode рассчитывается как $16\cdot\text{segment}+\text{offset}$.
Например, \texttt{1000:0000} обозначает физический \texttt{0x10000}.
Инструкции адресует \texttt{CS:IP}; \texttt{DS} указывает данные программы,
\texttt{SS:SP} --- её стек. В начале основного NASM-кода \texttt{DS=CS},
\texttt{SS=7000h}, \texttt{SP=F000h}. Глобальные данные и стек поэтому
не обязаны занимать одно место. Это учебная конфигурация, не модель
адресного пространства процесса современной ОС.

Чтобы вывести \texttt{S01  grade 07}, сначала используем готовый сервис
прошивки. В процедуре \texttt{bios\_text} символы читаются из строки
\texttt{DS:SI} и передаются в BIOS:
\begin{lstlisting}[basicstyle=\ttfamily\small]
lodsb                    ; AL = next character
test al, al
jz .done
mov ah, 0x0e             ; BIOS teletype
mov bx, 0x0007
int 0x10
\end{lstlisting}
\texttt{INT 10h} --- \emph{программная} инструкция обращения к BIOS;
это не аппаратное уведомление о событии. Цифра на экране --- ASCII-символ
\texttt{'7'}, а оценка 7 станет числовым значением данных лишь на следующей
стадии. Диагностический COM1 отдельно выдаёт \texttt{BOOT STAGE=1} и
\texttt{READY}, чтобы тест наблюдал действительный запуск без распознавания
изображения.

\paragraph{Версия 2: массив оценок и прямой вывод.}
Массив содержит 16 строк по 4 однобайтовых поля: $16\cdot4=64$ байта.
Коды 0--10 означают выставленные оценки, а \texttt{FF} означает
«не выставлена»; нуль \emph{не} может служить признаком пропуска.
Смещение выбранного элемента от начала массива равно $4i+j$.
В NASM это резервируется и инициализируется как
\texttt{grades: times 64 db 0xff}. Пояснение на C:
\begin{lstlisting}[language=C,basicstyle=\ttfamily\small]
uint8_t grade[16][4]; /* algorithm only */
\end{lstlisting}
Учебное значение \texttt{grade[1][2]} имеет смещение 6, но полный адрес
определяется сегментом \texttt{DS} и реальным размещением метки
\texttt{grades}. Адрес, байт значения и вывод на экран --- три разные вещи.

После знакомства с VGA текстовый вывод делаем самостоятельно: буфер
устройства начинается с физического \texttt{0xB8000}, то есть \texttt{B800:0000}.
На одну позицию приходятся код символа и байт атрибута. Строка экрана имеет
$80\cdot2=160$ байтов; адрес позиции внутри видеосегмента равен
$160\cdot\text{row}+2\cdot\text{col}$.
\begin{lstlisting}[basicstyle=\ttfamily\small]
mov ax, 0xb800
mov es, ax
; DI already contains the cell offset
lodsb                    ; AL = character at DS:SI
mov ah, 0x0f             ; color attribute
stosw                    ; AX -> ES:DI
\end{lstlisting}
Это ключевые строки реальной функции \texttt{vga\_text}.
Процедура \texttt{redraw} перебирает оценки в RAM и создаёт
\emph{отображение} в VGA; дисплей не становится долговременным хранилищем.
Для современной видеокарты потребовались бы иные драйверы и, возможно,
отдельная VRAM.

\paragraph{Версия 3: изменяем оценку с клавиатуры.}
Через BIOS \texttt{INT 16h} сначала спрашиваем о готовом событии
(\texttt{AH=01h}), затем читаем его (\texttt{AH=00h}):
\begin{lstlisting}[basicstyle=\ttfamily\small]
mov ah, 0x01
int 0x16
jz .loop
xor ah, ah
int 0x16
mov al, ah              ; AL = scan code for handle_scan
call handle_scan
\end{lstlisting}
PS/2-контроллер сообщает BIOS о клавише; BIOS предоставляет программе
высокоуровневый сервис. Клавиша \texttt{A} даёт оценку 10, цифры 0--9
задают соответствующее значение, Backspace записывает \texttt{FF},
стрелки выбирают работу и студента. Наш \texttt{handle\_scan} переводит
скан-код в действие. ASCII, скан-код и числовая оценка не тождественны.
Контроллер и CPU обмениваются через аппаратный интерфейс с адресом/портом,
данными и управляющей информацией; современные последовательные соединения
могут упаковывать это в сообщения по одним линиям. BIOS скрывал порты
\texttt{64h} и \texttt{60h} от первых версий программы.

\paragraph{Версия 4: собственное аппаратное прерывание.}
Теперь обработку PS/2 IRQ1 берёт наш код, поэтому BIOS \texttt{INT 16h}
больше не используется: его буфер клавиш уже не пополняется стандартным
обработчиком. В real mode таблица векторов \texttt{IVT} начинается с
\texttt{0000:0000}. Четыре байта по смещению $9\cdot4=36$ задают IP и CS
обработчика IRQ1 в совместимом PIC. При изменении IVT прерывания временно
выключаем через \texttt{cli}, затем включаем через \texttt{sti}.
После сохранения регистров короткий обработчик считывает скан-код
из порта \texttt{60h}, кладёт его в кольцевую очередь и подтверждает
PIC (EOI); возврат выполняет \texttt{iret}. Упрощённая схема кода:
\begin{lstlisting}[basicstyle=\ttfamily\small]
irq1:
    push ax
    push ds
    mov ax, cs
    mov ds, ax
    in al, 0x60
    ; enqueue the byte when there is room
    mov al, 0x20
    out 0x20, al
    pop ds
    pop ax
    iret
\end{lstlisting}
Реальный исходник также проверяет бит готовности, сохраняет \texttt{BX}
и содержит полный код очереди. Обработчик не перерисовывает экран:
основной цикл вызывает \texttt{dequeue}, затем \texttt{handle\_scan}.
Буфер из 64 байтов хранит не более 63 событий, потому что одна позиция
разделяет состояния «пусто» и «полно». Основной цикл всё ещё опрашивает
\emph{очередь}, но не PS/2-порт. Внешний IRQ1, намеренное \texttt{INT 16h}
и исключение от CPU имеют разные причины, хотя используют механизм
передачи управления.

\paragraph{Версия 5: статистика, регистры и память.}
Среднее вычисляем только из выставленных оценок. Если количество $n=0$,
показываем \texttt{--.--}, иначе в сотых:
\[
\left\lfloor\frac{100\cdot\mathrm{sum}+\lfloor n/2\rfloor}{n}\right\rfloor.
\]
Например, 0 и 10 дают 5.00; 0, 1 и 1 дают 0.67. Такое целочисленное
вычисление не требует библиотеки float. В NASM четыре поля читаются
из \texttt{DS:grades+BX}, байт \texttt{FF} пропускается, число 0
учитывается. Цикл из диагностики COM1:
\begin{lstlisting}[basicstyle=\ttfamily\small]
mov cx, WORKS
xor dx, dx             ; DH=sum, DL=count
.sum:
    mov al, [grades+bx]
    cmp al, UNSET
    je .skip
    inc dl
    add dh, al
.skip:
    inc bx
    loop .sum
\end{lstlisting}
Регистры держат текущую сумму; \texttt{CS:IP} выбирает байты машинных
инструкций, а \texttt{DS:grades+BX} --- байты данных. В строгой гарвардской
машине эти адресные пространства были бы разделены, в выбранном x86 PC
код и данные разделяют обычную RAM. Современные CPU могут иметь разные
L1I и L1D; QEMU-тест проверяет результат, но не измеряет настоящие
кэш-промахи. Последовательные 64 байта оценок демонстрируют локальность:
при условной 64-байтовой кэш-линии массив может занимать одну линию,
если выровнен, иначе две.

Стек используется вызовами, прерываниями и сохранением регистров.
В \texttt{entry} заранее устанавливаются \texttt{SS:SP}; данные оценок
лежат в отдельном статическом резерве. Для фиксированной группы
аллокатор кучи не требуется. Если разрешить динамическое число студентов,
понадобилось бы управление свободными блоками и временем жизни.

\paragraph{Версия 6: долговременное хранение.}
RAM энергозависима, поэтому оценки пишутся в отдельный диск.
Загрузочная программа на виртуальной дискете (диск BIOS 00h), оценки
в секторе LBA 1 первого HDD (диск BIOS 80h). BIOS \texttt{INT 13h}
с расширенным адресным пакетом DAP предоставляет чтение \texttt{AH=42h}
и запись \texttt{AH=43h}; флаг CF сообщает об ошибке. DAP задаёт
число секторов, физический адрес буфера \texttt{1000:sector} и номер LBA.
Здесь нет файловой системы: сектор не имеет имени файла. Как именно
BIOS взаимодействует с контроллером, этот вызов скрывает; при прямом
PIO CPU переносил бы слова сам, а при DMA контроллер передавал бы блок
после настройки. Ни DMA, ни собственный PIO-драйвер в действующей версии
журнала не реализованы.

\begin{center}\small
\begin{tabular}{|l|p{0.61\linewidth}|}\hline
Смещение & Формат одного сектора (512 байтов) \\ \hline
0--3 & ASCII \texttt{GR16} \\
4--7 & версия 1, 16 студентов, 4 работы, резерв 0 \\
8--71 & 64 оценки 0--10 или \texttt{FF} \\
72--73 & 16-битная сумма байтов 0--71, little-endian \\
74--511 & нули \\ \hline
\end{tabular}\end{center}
Новый формат \texttt{GR16} не смешиваем с историческим форматом
\texttt{GRD4} от x86-64-версии. Контрольная сумма обнаруживает некоторые
случайные повреждения, не защищает от преднамеренной подмены.
\texttt{save\_grades} зануляет 512-байтовый буфер, записывает поля по
фиксированным смещениям, считает сумму и вызывает \texttt{disk\_io}.
\texttt{load\_grades} проверяет диск, сигнатуру, версию, размеры,
сумму и все значения, только затем копирует таблицу в RAM. При ошибках
\texttt{EMPTY}, \texttt{NO\_DISK}, \texttt{BAD\_FORMAT},
\texttt{UNSUPPORTED\_VERSION}, \texttt{BAD\_CHECKSUM}, \texttt{BAD\_GRADE}
текущая таблица остаётся неизменной. Одна перезаписываемая копия сектора
не гарантирует атомарности при внезапной потере питания.

\paragraph{Машина в целом и переход к ОС.}
Плата и контроллеры соединяют CPU, RAM, накопитель, экран и клавиатуру.
В современном ПК для периферии используют USB, SATA и PCIe; NVMe ---
протокол хранения, а M.2 --- разъём/форм-фактор. Дискретный GPU
с собственной VRAM ориентирован на много подходящих параллельных операций;
наш текстовый VGA такой функциональности не реализует. HDD хранит данные
на вращающихся пластинах, SSD --- во Flash с контроллером. Их задержка
и пропускная способность различаются, хотя наш эмулируемый диск служит
лишь проверкой поведения, не скорости. BIOS смог помочь при запуске,
вводе и диске, но организации процессов, защиты памяти и файлов он
журналу не предоставил. Этим будет заниматься операционная система
в следующей лекции.

\paragraph{Самопроверка.}
Что произойдёт до \texttt{1000:0000}? Чем вызов \texttt{INT 16h}
отличается от IRQ1? Почему 0 --- оценка, а \texttt{FF} --- пропуск?
Почему \texttt{B800:0000} не равно адресу массива \texttt{grades}?
Что после перезапуска остаётся на дискете, а что на отдельном HDD?
Почему версия без гостевой ОС всё же может вызвать BIOS \texttt{INT 13h}?

Исходники действующего журнала: \path{demos/lecture04-journal/}.
Команда \texttt{make lecture4-journal-check} собирает шесть образов в Docker
и проверяет их запуск в QEMU; \texttt{make lecture4-journal-run}
показывает последнюю версию в браузере через noVNC по адресу
\url{http://localhost:6080/vnc.html}. Архивные 64-битные исходники
и прежняя презентация сохранены в \path{legacy/} для сравнения, но
не являются частью этого учебного маршрута.
\endgroup


\newpage
\section{Семинар №4. Память, адреса и представление данных.}

\subsection*{Цель занятия}
Закрепить представление о памяти как о последовательности байтов и связать это представление с простыми программами на C.

\subsection*{Обязательные задания}
\begin{itemize}
    \item Определить размер базовых типов C с помощью \texttt{sizeof}.
    \item Нарисовать размещение массива \texttt{char}, \texttt{int} и \texttt{long} в памяти.
    \item По адресу начала массива и размеру элемента вычислять адрес элемента с заданным индексом.
    \item Разобрать, чем отличаются значение переменной и адрес переменной.
    \item Составить таблицу: тип данных, размер, пример значения, возможное машинное представление.
\end{itemize}

\subsection*{Усложнённые задания}
\begin{itemize}
    \item Посмотреть байты целого числа в памяти и объяснить порядок байт.
    \item Разобрать размещение полей простой структуры C.
    \item Найти padding и выравнивание в структуре с полями разных типов.
    \item Сравнить содержимое памяти до и после изменения значения переменной.
\end{itemize}

\subsection*{Ключевые технологии и понятия}
\begin{itemize}
    \item \texttt{gcc} для сборки демонстрационных программ.
    \item \texttt{printf}, \texttt{sizeof}, оператор \texttt{\&} в C.
    \item \texttt{gdb}, \texttt{x/xb}, \texttt{x/xw} как дополнительные инструменты просмотра памяти.
\end{itemize}

\newpage
\chapter{Тема №4. Операционная система и запуск программы.}
\section{Лекция №5. Операционная система и запуск программы.}

\subsection*{Цель занятия}
Объяснить, зачем нужна операционная система и что происходит между запуском исполняемого файла и выполнением инструкций процессором.

\subsection*{Обязательные темы}
\begin{itemize}
    \item Задачи операционной системы: процессы, память, файлы, устройства, права доступа.
    \item Программа как файл на диске и процесс как выполняющийся экземпляр программы.
    \item Загрузка программы в память и передача управления начальной точке исполнения.
    \item Аргументы командной строки и переменные окружения.
    \item Стандартные потоки: \texttt{stdin}, \texttt{stdout}, \texttt{stderr}.
    \item Код завершения программы.
    \item Права доступа к файлам и исполняемый бит.
    \item Отличия Windows, Linux и macOS с точки зрения программиста.
\end{itemize}

\subsection*{Усложнённые темы}
\begin{itemize}
    \item Формат исполняемых файлов ELF.
    \item Системные вызовы как интерфейс между программой и ОС.
    \item Пользовательский режим и режим ядра.
    \item Изоляция процессов и виртуальное адресное пространство.
    \item Динамическая загрузка библиотек на уровне общей идеи.
\end{itemize}

\subsection*{Ключевые технологии и понятия}
\begin{itemize}
    \item Linux process model.
    \item ELF, loader, syscall.
    \item \texttt{file}, \texttt{chmod}, \texttt{./program}, \texttt{echo \$?}.
    \item \texttt{ps}, \texttt{top}, \texttt{strace} как дополнительные инструменты наблюдения.
\end{itemize}

\newpage
\section{Семинар №5. Процессы, файлы и стандартные потоки в Linux.}

\subsection*{Цель занятия}
Научиться запускать программы в Linux, понимать базовые ошибки запуска и пользоваться стандартными потоками ввода-вывода.

\subsection*{Обязательные задания}
\begin{itemize}
    \item Запустить готовую программу из текущего каталога через \texttt{./program}.
    \item Проверить тип файла командой \texttt{file}.
    \item Сделать файл исполняемым с помощью \texttt{chmod +x}.
    \item Передать программе аргументы командной строки.
    \item Перенаправить вывод в файл с помощью \texttt{>}.
    \item Передать файл на вход программе с помощью \texttt{<}.
    \item Соединить две команды через пайп \texttt{|}.
    \item Проверить код завершения последней команды через \texttt{echo \$?}.
\end{itemize}

\subsection*{Усложнённые задания}
\begin{itemize}
    \item Найти запущенный процесс через \texttt{ps}.
    \item Посмотреть системные вызовы простой программы через \texttt{strace}.
    \item Диагностировать ошибки \texttt{permission denied}, \texttt{command not found}, \texttt{no such file or directory}.
    \item Сравнить обычный вывод и поток ошибок при перенаправлении.
\end{itemize}

\subsection*{Ключевые технологии и понятия}
\begin{itemize}
    \item \texttt{stdin}, \texttt{stdout}, \texttt{stderr}.
    \item Exit code.
    \item \texttt{chmod}, \texttt{ps}, \texttt{strace}.
\end{itemize}

\newpage
\chapter{Тема №5. Linux как рабочая среда программиста.}
\section{Лекция №6. Терминал, shell и файловая система Linux.}

\subsection*{Цель занятия}
Дать минимальную модель Linux-среды, достаточную для выполнения учебных заданий, сборки программ и самостоятельной диагностики типовых проблем.

\subsection*{Обязательные темы}
\begin{itemize}
    \item Гипервизор, виртуальная машина, WSL2 и Docker: зачем нужны изолированные среды.
    \item Терминал и shell: кто читает команду пользователя.
    \item Файловая система Linux: корень \texttt{/}, домашний каталог, текущий каталог.
    \item Абсолютные и относительные пути.
    \item Каталоги \texttt{/home}, \texttt{/bin}, \texttt{/usr}, \texttt{/tmp} на уровне пользователя.
    \item Переменная окружения \texttt{PATH}.
    \item Работа с документацией через \texttt{man} и \texttt{--help}.
\end{itemize}

\subsection*{Усложнённые темы}
\begin{itemize}
    \item Пользователи, группы и права доступа.
    \item Символические ссылки.
    \item Устройство \texttt{/proc} как окна в состояние системы.
    \item Основы shell-скриптов: переменные, последовательность команд, код завершения.
    \item Отличие контейнера от виртуальной машины.
\end{itemize}

\subsection*{Ключевые технологии и понятия}
\begin{itemize}
    \item Ubuntu LTS как базовая учебная Linux-среда.
    \item Bash.
    \item VM, WSL2, Docker.
    \item \texttt{PATH}, current working directory, home directory.
\end{itemize}

\newpage
\section{Семинар №6. Навигация, поиск и работа с текстом.}

\subsection*{Цель занятия}
Отработать минимальный набор команд Linux, который нужен для работы с исходным кодом, файлами заданий и выводом программ.

\subsection*{Обязательные задания}
\begin{itemize}
    \item Узнать текущий каталог командой \texttt{pwd}.
    \item Перейти в другой каталог командой \texttt{cd}.
    \item Посмотреть содержимое каталога командой \texttt{ls}.
    \item Создать каталог и файл командами \texttt{mkdir} и \texttt{touch}.
    \item Скопировать, переместить и удалить файлы командами \texttt{cp}, \texttt{mv}, \texttt{rm}.
    \item Просмотреть текстовый файл с помощью \texttt{less}.
    \item Найти строки в файле с помощью \texttt{grep}.
    \item Найти файл по имени с помощью \texttt{find}.
    \item Посчитать строки, слова или байты с помощью \texttt{wc}.
\end{itemize}

\subsection*{Усложнённые задания}
\begin{itemize}
    \item Построить простой пайплайн из нескольких команд.
    \item Использовать регулярное выражение в \texttt{grep}.
    \item Найти все файлы с расширением \texttt{.c} или \texttt{.asm} в дереве каталогов.
    \item Написать небольшой shell-скрипт для запуска серии команд.
    \item Использовать \texttt{xargs} для обработки списка файлов.
\end{itemize}

\subsection*{Ключевые технологии и понятия}
\begin{itemize}
    \item \texttt{pwd}, \texttt{ls}, \texttt{cd}, \texttt{mkdir}, \texttt{touch}.
    \item \texttt{cp}, \texttt{mv}, \texttt{rm}.
    \item \texttt{less}, \texttt{grep}, \texttt{find}, \texttt{wc}, \texttt{head}, \texttt{tail}.
    \item Redirection, pipe, glob pattern.
\end{itemize}

\newpage
\chapter{Тема №6. Сборка программ и инструменты разработки.}
\section{Лекция №7. От исходного кода к исполняемому файлу.}

\subsection*{Цель занятия}
Показать путь от исходного текста программы до исполняемого файла и подготовить студентов к сборке программ на C и NASM.

\subsection*{Обязательные темы}
\begin{itemize}
    \item Исходный файл, объектный файл, исполняемый файл.
    \item Этапы сборки C-программы: препроцессор, компилятор, ассемблер, линковщик.
    \item Заголовочные файлы и единицы трансляции.
    \item Что делает \texttt{gcc} при обычном запуске.
    \item Что делает \texttt{nasm} при сборке ассемблерного файла.
    \item Зачем нужен линковщик.
    \item Зачем нужен \texttt{make} и файл \texttt{Makefile}.
    \item Типовые ошибки: ошибка компиляции, ошибка ассемблирования, ошибка линковки.
\end{itemize}

\subsection*{Усложнённые темы}
\begin{itemize}
    \item Секции исполняемого файла: \texttt{.text}, \texttt{.data}, \texttt{.bss}.
    \item Символы и таблица символов.
    \item Статическая и динамическая линковка на уровне общей идеи.
    \item ABI и соглашение о вызовах как мост между C и ассемблером.
    \item Отладочная информация и флаг \texttt{-g}.
\end{itemize}

\subsection*{Ключевые технологии и понятия}
\begin{itemize}
    \item \texttt{gcc}, \texttt{nasm}, \texttt{ld}, \texttt{make}.
    \item \texttt{Makefile}.
    \item \texttt{objdump}, \texttt{readelf}, \texttt{nm}.
    \item \texttt{gdb} как будущий инструмент отладки.
\end{itemize}

\newpage
\section{Семинар №7. Сборка C- и NASM-программ.}

\subsection*{Цель занятия}
Научиться собирать простые программы вручную и через \texttt{Makefile}, а также отличать ошибки компиляции от ошибок линковки.

\subsection*{Обязательные задания}
\begin{itemize}
    \item Собрать программу \texttt{hello.c} командой \texttt{gcc hello.c -o hello}.
    \item Запустить получившийся исполняемый файл.
    \item Разделить C-программу на два исходных файла и собрать их вместе.
    \item Получить объектный файл с помощью \texttt{gcc -c}.
    \item Написать минимальный \texttt{Makefile} с целями \texttt{all} и \texttt{clean}.
    \item Собрать простую NASM-программу в объектный файл.
    \item Слинковать объектный файл в исполняемый файл.
\end{itemize}

\subsection*{Усложнённые задания}
\begin{itemize}
    \item Посмотреть секции исполняемого файла через \texttt{readelf -S}.
    \item Посмотреть дизассемблирование через \texttt{objdump -d}.
    \item Найти символы через \texttt{nm}.
    \item Запустить программу под \texttt{gdb} и поставить точку останова на \texttt{main}.
    \item Специально вызвать ошибку линковки и объяснить её текст.
\end{itemize}

\subsection*{Ключевые технологии и понятия}
\begin{itemize}
    \item \texttt{gcc -c}, \texttt{gcc -o}.
    \item \texttt{nasm -f elf64}.
    \item \texttt{make}, \texttt{Makefile}, цель сборки.
    \item Object file, executable file, linker error.
\end{itemize}

\newpage
\chapter{Тема №7. Основы Git и командной разработки.}
\section{Лекция №8. Система контроля версий Git.}

\subsection*{Цель занятия}
Объяснить, зачем программисту система контроля версий, и дать модель Git, достаточную для учебных проектов и домашних заданий.

\subsection*{Обязательные темы}
\begin{itemize}
    \item Зачем нужна система контроля версий.
    \item Репозиторий, рабочая директория, индекс, коммит.
    \item История изменений и идентификатор коммита.
    \item Локальный и удалённый репозиторий.
    \item Ветка как линия разработки.
    \item Что такое конфликт и почему он возникает.
    \item Минимальная дисциплина работы: маленькие коммиты, понятные сообщения, не коммитить артефакты сборки.
\end{itemize}

\subsection*{Усложнённые темы}
\begin{itemize}
    \item Merge и rebase на уровне идеи.
    \item Pull request / merge request как способ проверки изменений.
    \item \texttt{.gitignore} и исключение временных файлов.
    \item Работа с несколькими ветками.
    \item Как читать \texttt{git diff} перед коммитом.
\end{itemize}

\subsection*{Ключевые технологии и понятия}
\begin{itemize}
    \item Git.
    \item GitHub или GitLab как удалённый хостинг репозиториев.
    \item Commit, branch, remote, conflict.
    \item \texttt{.gitignore}.
\end{itemize}

\newpage
\section{Семинар №8. Минимальный рабочий цикл Git.}

\subsection*{Цель занятия}
Отработать базовый цикл работы с репозиторием: изменить файл, проверить изменения, сделать коммит и отправить его в удалённый репозиторий.

\subsection*{Обязательные задания}
\begin{itemize}
    \item Создать локальный репозиторий командой \texttt{git init}.
    \item Проверить состояние рабочей директории командой \texttt{git status}.
    \item Добавить файл в индекс командой \texttt{git add}.
    \item Сделать коммит командой \texttt{git commit}.
    \item Посмотреть историю через \texttt{git log}.
    \item Посмотреть изменения через \texttt{git diff}.
    \item Склонировать удалённый репозиторий командой \texttt{git clone}.
    \item Получить и отправить изменения через \texttt{git pull} и \texttt{git push}.
\end{itemize}

\subsection*{Усложнённые задания}
\begin{itemize}
    \item Создать ветку и переключиться на неё с помощью \texttt{git branch} и \texttt{git switch}.
    \item Слить ветку обратно в основную.
    \item Создать учебный конфликт и разрешить его вручную.
    \item Написать \texttt{.gitignore} для проекта на C.
    \item Проверить изменения перед коммитом и убрать из коммита лишние файлы.
\end{itemize}

\subsection*{Ключевые технологии и понятия}
\begin{itemize}
    \item \texttt{git init}, \texttt{git status}, \texttt{git add}, \texttt{git commit}.
    \item \texttt{git log}, \texttt{git diff}.
    \item \texttt{git clone}, \texttt{git pull}, \texttt{git push}.
    \item \texttt{git branch}, \texttt{git switch}, \texttt{git merge}.
\end{itemize}
